\section{Discussion}
\label{sec:discussion}

Everything below is read off a \emph{single} two-model,
five-benchmark panel. We state the takeaways as what this case
study illustrates, not as comprehensive findings; their generality
is untested.

\paragraph{Why a self-audit chronology must be published.} F3b
(TruthfulQA MC2 correct-first ordering) and F3c (ToxiGen top-$k$
truncation) had already produced plausible publishable numbers
before per-cell scorer-output inspection exposed them. Any
assurance-grade benchmark audit should therefore publish a
self-audit chronology, not only cleaned final numbers. Box~1 gives
the minimal fields; Appendix~\ref{app:bugs} is our worked instance.
Without it, repair silently absorbed into a clean narrative is
indistinguishable from a pipeline where the bugs were never caught.
Appendix~\ref{app:verdict-mapping} maps the raw
\texttt{cell\_validity.csv} verdicts to Table~\ref{tab:eligibility},
making the eligibility counts mechanical.

\begin{tcolorbox}[
  title={Box~1: Minimal self-audit chronology template},
  fonttitle=\bfseries\small, fontupper=\footnotesize,
  colback=black!3, colframe=black!55]
\emph{Per failure encountered}, record: \textbf{(1)} a stable bug
ID and its failure class (e.g.\ F1--F5 or a project-local code);
\textbf{(2)} discovery date and \emph{how} it was found (which
inspection or test surfaced it---not just ``code review'');
\textbf{(3)} affected cells (every model\,$\times$\,benchmark the
defect touched); \textbf{(4)} the symptom \emph{as it appeared in
the reported numbers} (what a reader would have seen and believed);
\textbf{(5)} provenance---\emph{inherited} from the initial
pipeline vs.\ \emph{introduced during repair}; \textbf{(6)} the fix
(the specific code change) \emph{and} a regression test that fails
on the pre-fix code; \textbf{(7)} the before/after delta on the
headline statistic for each affected cell; \textbf{(8)} residual
status: \textsc{fixed}, \textsc{disclosed-unfixed}, or
\textsc{out-of-scope}, with the reason. \emph{Per chronology as a
whole}, also state: total entries, how many were repair-introduced,
and which remain unfixed and why. An entry missing field (4) or (7)
is not auditable: the reader cannot tell what the wrong number was
or how much it moved.
\end{tcolorbox}

\paragraph{Why we report four buckets, not one verdict.} The four
non-confirmatory statuses in \S\ref{sec:results} separate different
consumer actions: \emph{ineligible} calls for scope revision,
\emph{scorer-unvalidated} for scorer validation, \emph{failed
G2--G4} for threshold or evidence revision, and \emph{exploratory}
for disclosed, non-confirmatory release. Collapsing these into a
single ``inconclusive'' verdict would hide the reasoning step that a
governance reviewer most needs.

\paragraph{Zero confirmatory cells is procedural, not an empirical
null.} The zero confirmatory verdict is not a claim about the five
benchmarks. Under \S\ref{par:gate}, even a cell whose CI is wholly
above 1 remains \emph{exploratory} while G5 or G6 is
\textsc{proposed}; the two TruthfulQA cells would upgrade only after
those gates are implemented. The contribution is the gate and
taxonomy, not a benchmark-validity verdict. \emph{No
benchmark-validity claim is issued for any cell.}
